% answers_short.tex —— 2013 年·北京卷·理科 · 简短答案速查
% 由 tools/make_short_answers.py 自 answers.tex 的 \gksolhead/\gkanswer 提取，
% 并与 2013/2013_数学真题_北京理_答案要点.md 逐题比对后生成，勿手改（改源 MD 后重跑）。

\begin{gknotice}
\textbf{使用说明}：本节只给最终答案便于快速核对，完整推导见解析版。答案为编写者独立推导，\textbf{非官方标准答案}。
\end{gknotice}

\gkpart{选择题}
\begin{center}\begin{tabular}{ccccccccc}
\toprule
题号 & 1 & 2 & 3 & 4 & 5 & 6 & 7 & 8 \\
\midrule
答案 & B & D & A & C & D & B & C & C \\
\bottomrule
\end{tabular}\end{center}

\gkpart{填空题}
\begin{itemize}[leftmargin=2.6em]
  \item[\textbf{9.}] $1$。点的直角坐标为 $(\sqrt3,1)$，直线为 $y=2$，距离为 $|1-2|=1$。
  \item[\textbf{10.}] $q=2$，$S_n=2^{n+1}-2$。由 $a_3+a_5=q(a_2+a_4)$ 得 $q=2$；再由 $a_2+a_4=10a_1=20$ 得 $a_1=2$。
  \item[\textbf{11.}] $PD=\dfrac95$，$AB=4$。切割线定理给出 $PA^2=PD\cdot PB$，设 $PD=9k$，$DB=16k$，得 $k=\dfrac15$；又 $PA\perp AB$，故 $AB=\sqrt{PB^2-PA^2}=4$。
  \item[\textbf{12.}] $96$。连号的 $2$ 张有 $4$ 种选法，与其余 $3$ 张一起分给 $4$ 人有 $4!$ 种，共 $4\times4!=96$ 种。
  \item[\textbf{13.}] $4$。取网格坐标得 $\boldsymbol a=(-1,1)$，$\boldsymbol b=(6,2)$，$\boldsymbol c=(-1,-3)$，解得 $\lambda=-2$，$\mu=-\dfrac12$，故 $\dfrac{\lambda}{\mu}=4$。
  \item[\textbf{14.}] $\dfrac{2\sqrt5}5$。建系得 $P(2t,2-t,2-2t)$，到直线 $CC_1$ 的距离为 $\sqrt{5t^2-8t+4}$，在 $t=\dfrac45$ 处取最小值 $\dfrac{2\sqrt5}5$。
\end{itemize}

\gkpart{解答题}
\begin{itemize}[leftmargin=2.6em]
  \item[\textbf{15.}] (1) $\dfrac{\sqrt6}3$。由正弦定理及 $\angle B=2\angle A$ 得 $3=2\sqrt6\cdot\dfrac{\sin A}{\sin2A}=\dfrac{\sqrt6}{\cos A}$；(2) $5$。$\cos B=\dfrac13$，$\sin B=\dfrac{2\sqrt2}3$，$\sin C=\sin(A+B)=\dfrac{5\sqrt3}9$，故 $c=\dfrac{a\sin C}{\sin A}=5$
  \item[\textbf{16.}] (1) $\dfrac2{13}$。$3$ 月 $1$ 日至 $13$ 日中等概率选一天，到达当日重度污染只有 $5$ 日与 $8$ 日两天；(2) $P(X=0)=\dfrac5{13}$，$P(X=1)=\dfrac4{13}$，$P(X=2)=\dfrac4{13}$，$E(X)=\dfrac{12}{13}$。逐日检查相邻两天中空气质量优良的天数即得；(3) 从 $5$ 日开始。三天 $220$，$160$，$40$ 的方差为 $5600$，大于其余各起始日对应的方差
  \item[\textbf{17.}] (1) 证明见解析。平面 $ABC\perp$ 平面 $AA_1C_1C$ 且交线为 $AC$，正方形中 $AA_1\perp AC$，故 $AA_1\perp$ 平面 $ABC$；(2) $\dfrac{16}{25}$。两半平面的法向量可取 $(0,4,3)$ 与 $(3,4,0)$，夹角余弦为 $\dfrac{16}{5\times5}=\dfrac{16}{25}$；(3) $\dfrac9{25}$。设 $D=B+t(C_1-B)$，由 $AD\cdot A_1B=9-25t=0$ 得 $t=\dfrac9{25}$，且 $0<\dfrac9{25}<1$，故点 $D$ 存在
  \item[\textbf{18.}] (1) $y=x-1$。$y'=\dfrac{1-\ln x}{x^2}$，在 $x=1$ 处斜率为 $1$，切线过 $(1,0)$；(2) 证明见解析。令 $h(x)=x^2-x-\ln x$，则 $h'(x)=\dfrac{(2x+1)(x-1)}x$，$h$ 在 $x=1$ 处取最小值 $0$，故 $x\neq1$ 时 $\dfrac{\ln x}{x}<x-1$
  \item[\textbf{19.}] (1) $\sqrt3$。由 $|OA|=|AB|$ 得 $A$ 的横坐标为 $1$，$|y|=\dfrac{\sqrt3}2$，菱形面积 $=|OB|\cdot|AC|/2=\sqrt3$；(2) 不可能。若 $OABC$ 为菱形，则 $OB$ 与 $AC$ 互相垂直平分；设 $B=(x_0,y_0)$，由 $A$，$C$ 都在 $W$ 上推出 $3y_0v_y=0$，与 $B$ 不是顶点矛盾
  \item[\textbf{20.}] (1) $d_1=1$，$d_2=1$，$d_3=3$，$d_4=3$。$A_n$ 依次为 $2,2,4,4$，而 $B_n$ 恒为 $1$；(2) 证明见解析。等差数列时不减，$A_n=a_n$，$B_n=a_{n+1}=a_n+d$；反之由 $B_n=A_n+d$ 与 $B_n\leqslant a_{n+1}$ 可推出 $a_{n+1}=A_n+d$，归纳得公差为 $d$；(3) 证明见解析。由 $B_1=1$ 知各项不小于 $1$；若某项 $\geqslant3$ 则 $B_n$ 与 $A_{n+1}$ 的关系矛盾，故各项只能是 $1$ 或 $2$；又每个 $B_n=1$ 都要求 $n$ 项之后出现 $1$，故有无穷多项为 $1$
\end{itemize}
